\section{Reduction of Spectra}

With longslit spectroscopy, to reduce parallactic angle
\index{slit!parallactic angle} \index{parallactic angle!slit}
\glossary{g-parallactic-angle} the slit is usually rotated to be along
a perpendicular line of azimuth (from the zenith to the horizon, and normal to the
horizon). This means the coordinate system for the slit is not simply the
declination of the object. However for a single star with a stated RA and
Dec that point on the slit corresponds to the actual RA and Dec in the sky.

The slit, grating, and other optics inside the spectrograph produce a long
\dhl{trace} \index{trace} onto the sensor. With several objects, these traces are referred to
as apertures. \index{spectra!aperture}. Any independent light path that produces a trace
is referred to as a \dhl{beam}. \index{beam!spectra} \index{spectra!beam} \gls{g-beam}.
You will see the term \dhl{band} - this will refer to the order of the spectra, or in
the case of the multispec format different process outputs for a single trace that include
sky noise and with the proper keywords the SNR of the trace.

The steps are:


\begin{itemize}
\addtolength{\itemsep}{-0.5\baselineskip}
   \item   Make an ``Observations'' directory for the night, and create a RawData directory, and
a \dhl{usw} directory.
   \item   Copy the template \dhl{pyraf.reduce} file into \dhl{usw}. 
   \item   Open the \dhl{pyraf.reduce} in the editor and start the commands/audit trail.
   \item   Properly archive the data.
   \item   Copy the raw data into the RawData directory.
   \item   Copy the RawData files into the PreAnalysis directory.\end{itemize}



Working in the PreAnalysis directory:

\begin{itemize}
\addtolength{\itemsep}{-0.5\baselineskip}
   \item   fix the filenames. Remove spaces, change +/- to p/m.
   \item   run fitserial to add a unique serial number.
   \item   update the headers as needed, remember DISPAXIS=1 is critical.
   \item   Alter the IMAGETYP keyword:
   \item   Make them all IMAGETYP=object,
   \item   then change the Zero* files to IMAGETYP=zero
   \item   trim the images. This reduces memory requirements, critical on a Mac.\\
  at this time the files should have a \dhl{t\_} prepended to the filename.
   \item   create the zeromaster.fits file.
   \item   ``zero subtract'' the zero from all the files but the Zero* files.\\
  at this time the files will have \dhl{z\_t\_} prepended to the filename.\end{itemize}



Now the fun begins.



Each target, we will call a sequence. Each sequence consists
of 3x target observations, 3x calibration lamp images (cals),
and 3x flat images. if the instrument has any flexure, flats
are necessary in each sequence.

The steps are roughly:

Interactively use apall for the middle of the science images. You can re-run apall
for the other two in non-interactive mode. (Yes, make sure the step is done right).

Next using the middle image as a reference, use apall,
non-interactive, to extract the matching trace pixels from the cal and flat images.

Use the \dhl{identify} step with the cal to get the WCS information for the
files in the sequence. 

Apply the WCS with \dhl{disport}. 

% HEREHEREHERE
